\documentclass[11pt,a4paper]{article}

\usepackage[T1]{fontenc}
\usepackage[utf8]{inputenc}
\usepackage{mathptmx}
\usepackage[expansion=false]{microtype}
\usepackage[margin=2.5cm]{geometry}
\usepackage{amsmath,amssymb}
\usepackage{booktabs}
\usepackage{array}
\usepackage{siunitx}
\usepackage{xcolor}
\usepackage{graphicx}
\usepackage{listings}
\usepackage{enumitem}
\usepackage{fancyhdr}
\usepackage{url}
\Urlmuskip=0mu plus 1mu
\usepackage[colorlinks=true,linkcolor=blue!50!black,citecolor=blue!50!black,urlcolor=blue!50!black]{hyperref}

\definecolor{codebg}{RGB}{245,245,245}
\definecolor{warncol}{RGB}{160,60,0}
\lstdefinestyle{dv}{
  backgroundcolor=\color{codebg},
  basicstyle=\ttfamily\small,
  commentstyle=\color{gray!70!black},
  morecomment=[l]{\#},
  breaklines=true,
  frame=single, framerule=0pt,
  xleftmargin=6pt, xrightmargin=6pt,
  aboveskip=8pt, belowskip=8pt
}

\newcommand{\code}[1]{\path{#1}}
\newcommand{\file}[1]{\path{#1}}
\newcommand{\cardline}[1]{\texttt{#1}}
\newcommand{\Interface}{\path{krc_davinciinterface.dvrc}}
\newcommand{\warn}{\textcolor{warncol}{\textbf{!}}\ }
\newcommand{\dif}{\mathrm{d}}

\pagestyle{fancy}
\fancyhf{}
\fancyhead[L]{\small\sffamily Extending the KRC--Davinci interface}
\fancyhead[R]{\small\thepage}
\renewcommand{\headrulewidth}{0.3pt}

\title{\textbf{Extending the KRC--Davinci Interface}\\[6pt]
\large Subsurface solar energy deposition, temperature-dependent emissivity,
adaptive meshing, and ice thermal conductivity}
\author{L.~Lange}
\date{}

\begin{document}
\maketitle
\thispagestyle{empty}

\begin{abstract}
\noindent
This note documents four additions made to the Davinci interface that drives KRC
(\Interface, function \code{krc()}, and its helper \code{krc_evalN1}, originally
written by S.~Piqueux): (i) distribution of absorbed solar energy within the
subsurface rather than at the surface alone, (ii) a temperature-dependent surface
emissivity, (iii) an adaptive first-layer thickness that keeps the near-surface mesh
consistent with (i), and (iv) temperature-dependent thermal conductivity and density
for water ice. The first two correspond to new KRC input cards (17 and 18); the
mesh adaptation is interface-side bookkeeping required by the first; the ice
properties feed the existing $T$-dependent conductivity card (12) through a
material-property lookup. Prescribing MCD/PCM atmospheric fluxes in KRC
(\code{UserPrescribed_Fluxes}, card 8/26) is a fifth, already-documented change:
it is only summarized here, with pointers to the dedicated design document.
\end{abstract}

\section{Context}

\Interface\ builds a KRC input deck from high-level physical parameters, sizes the
default vertical mesh (\code{krc_evalN1}), and manages external flux and albedo
files. It is the layer between physically meaningful arguments (thermal inertia,
albedo, orbital elements, \dots) and the fixed-format card deck that the KRC Fortran
executable reads (see \file{docs/flux_tables.md} and the two new-feature notes,
\file{docs/solar_penetration_in_ground.md} and
\file{docs/emissivity_temperature_dependant.md}, in the KRC repository). The four
modifications below extend both the KRC side (two new card types) and the interface
side (argument list of \code{krc()}, mesh sizing logic, and a material-property
helper).

\section{Subsurface solar energy deposition (RADGND / EFOLD\_RADGND)}
\label{sec:radgnd}

By default KRC deposits all absorbed insolation $S_0(1-A)$ at the surface, which
appears as a source term in the surface radiative balance. For optically
translucent materials (e.g.\ water ice), a fraction of that energy is instead
absorbed in depth. \code{RADGND} switches on a volumetric deposition following a
Beer--Lambert-type exponential profile with e-folding length
\code{EFOLD_RADGND} (in metres):
\begin{equation}
  \frac{\dif Q}{\dif z} \;\propto\; \exp\!\left(-\frac{z}{\ell}\right), \qquad
  \ell \equiv \texttt{EFOLD\_RADGND},
\end{equation}
and removes the corresponding term from the surface boundary condition, which then
reduces to $\varepsilon\sigma T_\mathrm{surf}^4 = -k\,\partial T/\partial z$ (no
$S_0(1-A)$ term at $z=0$). On the interface side this is exposed as two new
\code{krc()} arguments, translated into KRC card 17:

\begin{center}
\begin{tabular}{@{}p{2.7cm}p{5.3cm}p{5.8cm}@{}}
\toprule
Argument & Meaning & KRC card \\
\midrule
\code{RADGND} & enable subsurface deposition (boolean) & \cardline{17 1 \textless{}0/1\textgreater{} 'RADGND' /} \\
\code{EFOLD_RADGND} & e-folding length $\ell$ [m] & \cardline{17 2 \textless{}val\textgreater{} 'EFOLD\_RADGND' /} \\
\bottomrule
\end{tabular}
\end{center}

\warn A mesh sized only for surface thermal diffusion can be far coarser than
$\ell$; Section~\ref{sec:flay} describes the mesh-side consequence of enabling
\code{RADGND}. Validation: \file{Test_effectTIvssolarpene},
\file{Test_constantabslength_vs_cosilength} (incidence-angle scaling of the
absorption length), against the analytical reference in
\file{Validation_solarpenetration.dv}.

\section{Temperature-dependent surface emissivity (EmisT)}
\label{sec:emist}

\code{EmisT} replaces the constant surface emissivity by a third-order polynomial
in temperature, evaluated at each iteration of the surface energy balance:
\begin{equation}
  \varepsilon(T) = \mathtt{Emis0} + \mathtt{Emis1}\,X + \mathtt{Emis2}\,X^2 +
  \mathtt{Emis3}\,X^3, \qquad X = \frac{T - 220}{100},
\end{equation}
the same normalized-temperature basis already used by the $T$-dependent
conductivity and specific-heat cards (12/13). The result is clamped to
$[0,1]$; a non-positive value aborts the current case (KRC's blow-up handler), so
coefficient sets should be checked against the expected surface temperature range
before production use. Interface side, this is card 18:

\begin{center}
\begin{tabular}{@{}lll@{}}
\toprule
Argument & Meaning & KRC card \\
\midrule
\code{EmisT} & enable $T$-dependent emissivity (boolean) & \cardline{18 1 \textless{}0/1\textgreater{} 'EmisT' /} \\
\code{Emis0}--\code{Emis3} & polynomial coefficients & \cardline{18 2..5 \textless{}val\textgreater{} 'Emis0'..'Emis3' /} \\
\bottomrule
\end{tabular}
\end{center}

Coefficients are obtained by fitting $\varepsilon(T)$ derived from a Mie/Hapke
emissivity model for amorphous and crystalline water ice
(\file{Compute_Fit_Emissivity_4KRC.dvrc}, inputs
\file{emis_amorphous_ice.txt}, \file{emis_crystalline_ice.txt}); fit quality is
checked in \file{Figures/Figure_EmissIcevsT.pdf} against the tabulated true
emissivity for a range of grain radii and absorption coefficients. The combined
effect of \code{EmisT} and $T$-dependent conductivity on $T_\mathrm{surf}$ is
checked separately in \file{combined_effect_emissT_and_conductivity}.

\section{Adaptive first-layer thickness (FLAY)}
\label{sec:flay}

\code{krc_evalN1} sizes the default vertical mesh (first-layer thickness
\code{FLAY}, layer count \code{IC2}/\code{N1}) from the diurnal and seasonal skin
depths implied by the surface thermal inertia, density, and specific heat, for
three layering cases (homogeneous half-space, fixed-thickness two-layer medium, or
an exponential Hayne-type profile). When \code{RADGND} is active, the resulting
\code{FLAY} is additionally capped so that the first mesh element stays fine
compared to the solar e-folding length, guaranteeing that the volumetric heating
profile of Section~\ref{sec:radgnd} is numerically resolved:
\begin{equation}
  \text{if } \mathtt{RADGND=T} \text{ and } \mathtt{FLAY2} >
  \frac{\ell}{3\,\mathtt{RLAY}\cdot\mathtt{RLAY}\cdot\mathtt{DSD1\_m}\cdot(0.5+(1-\mathtt{RLAY}))},
  \quad \mathtt{FLAY2} \leftarrow \text{that bound},
\end{equation}
where \code{DSD1_m} is the diurnal skin depth of the top material and
\code{RLAY} the layer-thickness growth ratio. Without this cap, a mesh dimensioned
for surface thermal diffusion alone can be coarser than $\ell$, biasing the
in-depth heating profile. Validated in \file{Test_adaptativeFLAY} by comparing
fixed-\code{FLAY} and adaptive-\code{FLAY} runs at two thermal inertias (33 and 133
$\mathrm{J\,m^{-2}\,K^{-1}\,s^{-1/2}}$) and checking convergence of $T_\mathrm{surf}$.

\section{Temperature-dependent ice thermal properties}
\label{sec:icetprop}

The material-property helper \code{Mat_Prop} returns, for a given material, the
polynomial coefficients expected by KRC's $T$-dependent conductivity and specific
heat cards (12/13), on the same normalized basis $X=(T-220)/100$ as
Section~\ref{sec:emist}. The \code{"H2O"} entry encodes:
\begin{itemize}[nosep]
  \item thermal conductivity $k(T)$ after Slack (1980), valid
  \SIrange{60}{270}{\kelvin};
  \item density $\rho(T)$ after R\"ottger et al.\ (2012), valid
  \SIrange{10}{265}{\kelvin};
  \item specific heat $C_p(T)$ after Giauque \& Stout (1936), valid
  \SIrange{15}{273}{\kelvin}.
\end{itemize}
These defaults are derived from, and cross-checked against, an icy-regolith
conductivity model reproducing Ferrari \& Lucas (\file{Scripts_reproduceFerrarriandLucas})
and initially refit for standalone injection into KRC
(\file{Script_Fiticeproperties4KRC}). End-to-end sensitivity of KRC output to the
new coefficients (surface temperature and effective thermal inertia, as a function
of albedo and stratigraphy) is checked in \file{Test_newcondinKRC}.

\paragraph{Joint fit with emissivity.} Since the conductivity of icy regolith
itself depends on emissivity (radiative conductivity term), grain radius and
porosity, fitting conductivity and emissivity independently is only an
approximation. \file{Compute_Fit_PropIcyRegolith_4KRC.dvrc} replaces the two
separate fits above by a single, self-consistent one (Davinci function
\code{Compute_Fit_ThermophysicalProperties_IcyRegolith}): given a grain radius,
porosity, conductivity mixing model (\code{GJKRBB}/\code{GB}/\code{WBB}/\code{WGB})
and ice form (amorphous/crystalline), it returns the \code{Emis0}--\code{Emis3}
coefficients of Section~\ref{sec:emist} together with matching
\emph{Con0--Con3} and \emph{Sph0--Sph3} coefficients for card 12/13, recomputing
emissivity internally when it is temperature-dependent rather than assuming a
fixed default. This is the current, production version of this fit; the
effect of enabling each $T$-dependence (conductivity, emissivity, or both)
on $T_\mathrm{surf}$ is checked in \file{joint_fit_validation}.

\section{Summary of new \texttt{krc()} arguments}

\begin{center}
\begin{tabular}{@{}p{3.1cm}p{6cm}p{5.5cm}@{}}
\toprule
Argument & Purpose & KRC card \\
\midrule
\code{RADGND} & enable subsurface solar deposition & \cardline{17 1 \textless{}0/1\textgreater{} 'RADGND' /} \\
\code{EFOLD_RADGND} & solar absorption e-folding length [m] & \cardline{17 2 \textless{}val\textgreater{} 'EFOLD\_RADGND' /} \\
\code{EmisT} & enable $T$-dependent emissivity & \cardline{18 1 \textless{}0/1\textgreater{} 'EmisT' /} \\
\code{Emis0}--\code{Emis3} & $\varepsilon(T)$ polynomial coefficients & \cardline{18 2..5 \textless{}val\textgreater{} 'Emis0'..'Emis3' /} \\
\code{UserPrescribed_Fluxes} & prescribed VIS/IR/ATM flux table (MCD/PCM) & \cardline{8 26 0 '\textless{}file\textgreater{}' /} \\
\bottomrule
\end{tabular}
\end{center}

\code{UserPrescribed_Fluxes} is unchanged by the present work and is documented in
full (format, database construction, coupling tutorial) in the companion design
document \emph{Prescribing Mars PCM/MCD Atmospheric Fluxes in KRC}; it is listed
here only for completeness, since it shares the interface file and the same
argument-to-card mechanism as the four additions above.

\section{Validation trace}

Complete inputs, scripts and figures are archived alongside this note:
\code{KRC_repo_additions/docs/validation/} for the two KRC-side cards (RADGND,
EmisT --- Sections~\ref{sec:radgnd}--\ref{sec:emist}); adaptive FLAY validation
sits with it (\file{Test_adaptativeFLAY}, Section~\ref{sec:flay}, since it is
checked through the RADGND card). On the interface side,
\code{Davinci_interface/CouplingwithMCD_PCM/} covers \code{UserPrescribed_Fluxes},
and \code{Davinci_interface/Improve_IceProperties/} covers the ice thermal
properties of Section~\ref{sec:icetprop}, including the joint fit.
A short \file{README} in every folder of that archive states its content and how it
maps to the sections above.

\end{document}
